\documentclass[11pt]{amsart}

\usepackage[T1]{fontenc}
\usepackage{lmodern}
\usepackage{microtype}
\usepackage{amsmath,amssymb,amsthm}
\usepackage[hidelinks]{hyperref}

\newtheorem{theorem}{Theorem}
\newtheorem{corollary}[theorem]{Corollary}

\title[Four factorizations of $K_{35}$]{Four Factorizations of \texorpdfstring{$K_{35}$}{K\_35}
Realizing \texorpdfstring{$\beta\in\{1,2,4,6\}$}{beta = 1, 2, 4, 6} in the Hamilton--Waterloo
Problem \texorpdfstring{$\mathrm{HWP}(35;5,7)$}{HWP(35;5,7)}}
\date{}
\subjclass[2020]{Primary 05C70; Secondary 05B30, 05C51, 05C38}
\keywords{Hamilton--Waterloo problem, two-factorization, cycle system,
resolvable decomposition, prescribed automorphism, complete graph}

\hypersetup{
  pdftitle={Four factorizations of K_35 realizing beta = 1, 2, 4, 6 in the Hamilton-Waterloo problem HWP(35;5,7)}
}

\newcommand{\HWP}{\mathrm{HWP}}

\begin{document}

\begin{abstract}
We print four factorizations of $K_{35}$ into seventeen two-factors,
realizing $(\alpha,\beta)=(16,1)$, $(15,2)$, $(13,4)$ and $(11,6)$ in
the Hamilton--Waterloo problem $\HWP(35;5,7)$. Each object is given in
full as a list of cycles, and checking one is a finite computation on
that list. Burgess, Danziger and Traetta determined $\HWP(mn;m,n)$ for
odd $n\ge m\ge3$ up to an explicit list of possible exceptions, whose
clause at $(m,n)=(5,7)$ names the cells with $\beta\in\{1,2,4,6\}$ and
the quadruple $(\alpha,\beta)=(9,8)$. Section~1 paraphrases, from its
e-print, a theorem of Wang, Lu and Cao covering the fourteen pairs at
$v=35$ with $\beta\notin\{1,2,4,6\}$; granting that statement, which is
neither displayed nor verified here, these four factorizations account
for the remaining four of the eighteen pairs with $\alpha,\beta\ge0$ and
$\alpha+\beta=17$. No priority is claimed for the four cells.
\end{abstract}

\maketitle

\section{The problem and the two quoted theorems}

For $v\ge3$ and $3\le m,n\le v$, the Hamilton--Waterloo problem
$\HWP(v;m,n)$ asks for which pairs $(\alpha,\beta)$ the edge set of
$K_v$ decomposes into $\alpha$ two-factors all of whose cycles have
length $m$ and $\beta$ two-factors all of whose cycles have length $n$.
We write $(\alpha,\beta)\in\HWP(v;m,n)$ when such a decomposition
exists. For odd $v$ a two-factorization of $K_v$ has exactly
$(v-1)/2$ factors, so $\alpha+\beta=(v-1)/2$ is forced; and $m\mid v$
is forced when $\alpha\ge1$, $n\mid v$ when $\beta\ge1$. At
$(m,n,v)=(5,7,35)$ both divisibilities hold, so the only necessary
condition in force below is $\alpha+\beta=17$.

Burgess, Danziger and Traetta \cite{BDT17} proved, for odd $n\ge m\ge3$
and $v=mn$, that $(\alpha,\beta)\in\HWP(mn;m,n)$ if and only if
$\alpha,\beta\ge0$ and $\alpha+\beta=(mn-1)/2$, \emph{except possibly}
when
\[
  \beta\in\Bigl[1,\tfrac{n-3}{2}\Bigr]\cup
  \Bigl\{\tfrac{n+1}{2},\ \tfrac{n+5}{2}\Bigr\},
  \qquad\text{or}\qquad
  (m,\alpha)\in\{(3,2),(3,4)\},
\]
or $(m,n,\alpha,\beta)$ is one of seven listed quadruples. This
statement, and that of \cite{WLC} paraphrased below, are read off the
e-prints of those papers rather than from the typeset journal pages;
neither is reproduced verbatim here, and nothing below verifies either
of them.

At $(m,n)=(5,7)$ we have $v=35$ and $\alpha+\beta=17$. Here
$[1,(n-3)/2]=\{1,2\}$ and $\{(n+1)/2,(n+5)/2\}=\{4,6\}$, the clause
$(m,\alpha)\in\{(3,2),(3,4)\}$ is inert because $m=5$, and exactly one
of the seven listed quadruples has $(m,n)=(5,7)$, namely
$(5,7,9,8)$. So the exception clause of \cite{BDT17} names five cells at
$v=35$, those with $\beta\in\{1,2,4,6,8\}$.

On our reading of the e-print of \cite{WLC}, Wang, Lu and Cao prove
$(9,8)\in\HWP(35;5,7)$ outright, and their general theorem for odd
orders excepts, at $v=35$, only $\beta\in\{1,2,4,6\}$; that is, it
gives the fourteen pairs with $\beta\notin\{1,2,4,6\}$, the cell
$(9,8)$ among them. We use this only as quoted, in
Corollary~\ref{cor:complete}. The four factorizations below realize the
pairs with $\beta\in\{1,2,4,6\}$; whether those four pairs had been
realized before is not settled here.

\section{Two prescribed automorphisms}

Take $V(K_{35})=\{0,1,\dots,34\}$. We use two permutations of $V$, each
of prime order, and each acting on $V$ in blocks of consecutive
integers.

\emph{The $p=7$ lane.} Write $u=7c+x$ with $c\in\{0,\dots,4\}$ and
$x\in\mathbb{Z}_7$, and set
\[
  \rho(7c+x)=7c+\bigl((x+1)\bmod 7\bigr).
\]
Thus $\rho$ has order $7$ and fixes each of the five \emph{classes}
\[
  \{0,\dots,6\},\ \{7,\dots,13\},\ \dots,\ \{28,\dots,34\}
\]
setwise.

\emph{The $p=5$ lane.} Write $u=5j+y$ with $j\in\{0,\dots,6\}$ and
$y\in\mathbb{Z}_5$, and set
\[
  \sigma(5j+y)=5j+\bigl((y+1)\bmod 5\bigr),
\]
of order $5$, fixing each of the seven classes
\[
  \{0,\dots,4\},\ \{5,\dots,9\},\ \dots,\ \{30,\dots,34\}
\]
setwise.

Both actions on $E(K_{35})$ are semiregular, so all edge orbits have
size $p$. For $\rho$ there are $85$ orbits: fifteen \emph{intra-class}
orbits, the three circulants $\mathrm{Cay}(\mathbb{Z}_7,\{\pm d\})$,
$d=1,2,3$, inside each of the five classes, and seventy
\emph{inter-class} orbits, one for each of the ten pairs of classes and
each of the seven shifts; and $85\times7=595=\binom{35}{2}$. For
$\sigma$ there are $119$ orbits, $14$ intra-class and $105$
inter-class, and $119\times5=595$. An intra-class orbit is a $p$-cycle
on its class; an inter-class orbit is a perfect matching between its
two classes.

A two-factor of $K_{35}$ fixed setwise by the lane map is called
\emph{invariant}; it is a union of $35/p$ edge orbits and counts as one
factor. A two-factor $F$ whose images $F,\rho F,\dots,\rho^{p-1}F$
(respectively $F,\sigma F,\dots$) are pairwise edge-disjoint is called a
\emph{base} object; it stands for $p$ factors and uses up $35$ edge
orbits. This is the whole bookkeeping: a certificate is a short list of
objects whose expansion must have $17$ factors using each of the $595$
edges exactly once.

\section{The four cells}

\begin{theorem}\label{thm:main}
$(16,1)$, $(15,2)$, $(13,4)$ and $(11,6)$ all lie in
$\HWP(35;5,7)$.
\end{theorem}

\begin{proof}
Section~\ref{sec:certs} prints, for each of the four cells, a list of
objects in one of the two lanes. Expanding each base object into its
$p$ translates gives in every case $17$ two-factors; the verification
of Section~\ref{sec:check} confirms that each factor is a $C_5$- or
$C_7$-factor as declared, that the $17$ factors are pairwise
edge-disjoint with union all of $E(K_{35})$, and that the number of
$C_5$- and $C_7$-factors is the claimed pair $(\alpha,\beta)$.
\end{proof}

\begin{corollary}\label{cor:complete}
Suppose the theorem of \cite{WLC} gives the fourteen pairs
$(\alpha,\beta)$ with $\alpha,\beta\ge0$, $\alpha+\beta=17$ and
$\beta\notin\{1,2,4,6\}$, as its e-print states and as paraphrased in
Section~1. Then
\[
  \HWP(35;5,7)=\{(\alpha,\beta):\alpha,\beta\ge0,\ \alpha+\beta=17\};
\]
that is, all eighteen admissible pairs are realizable.
\end{corollary}

\begin{proof}
$\alpha+\beta=17$ is necessary because $35$ is odd. The fourteen pairs
with $\beta\notin\{1,2,4,6\}$ are the hypothesis, and the four
with $\beta\in\{1,2,4,6\}$ are Theorem~\ref{thm:main}.
\end{proof}

The hypothesis of Corollary~\ref{cor:complete} is quoted from an e-print
and is neither displayed nor verified here; only the four cells of
Theorem~\ref{thm:main} are established below.

\section{The certificates}\label{sec:certs}

Each object is printed as its decomposition into cycles: a $C_5$-factor
as seven $5$-cycles and a $C_7$-factor as five $7$-cycles, together
spanning $\{0,\dots,34\}$. An object marked \emph{base} stands for
itself and its $p-1$ translates under the lane map; an object marked
invariant stands for one factor and is fixed setwise by that map.

\subsection*{$(\alpha,\beta)=(16,1)$, lane $p=7$}
\smallskip\noindent $G_{1}$: a $C_7$-factor, $\rho$-invariant.
\begin{gather*}
  (0,2,4,6,1,3,5)(7,9,11,13,8,10,12) \\
  (14,16,18,20,15,17,19)(21,24,27,23,26,22,25) \\
  (28,31,34,30,33,29,32)
\end{gather*}
\smallskip\noindent $G_{2}$: a $C_5$-factor, \emph{base}.
\begin{gather*}
  (0,12,11,23,31)(1,22,18,28,33)(2,8,4,5,19) \\
  (3,10,13,26,16)(6,7,24,30,25)(9,15,27,21,32) \\
  (14,17,29,20,34)
\end{gather*}
\smallskip\noindent $G_{3}$: a $C_5$-factor, \emph{base}.
\begin{gather*}
  (0,22,2,5,34)(1,10,28,8,26)(3,13,21,30,29) \\
  (4,27,25,17,32)(6,14,15,24,18)(7,31,20,11,33) \\
  (9,19,12,16,23)
\end{gather*}
\smallskip\noindent $G_{4}$: a $C_5$-factor, $\rho$-invariant.
\begin{gather*}
  (0,16,11,27,30)(1,17,12,21,31)(2,18,13,22,32) \\
  (3,19,7,23,33)(4,20,8,24,34)(5,14,9,25,28) \\
  (6,15,10,26,29)
\end{gather*}
\smallskip\noindent $G_{5}$: a $C_5$-factor, $\rho$-invariant.
\begin{gather*}
  (0,18,10,31,24)(1,19,11,32,25)(2,20,12,33,26) \\
  (3,14,13,34,27)(4,15,7,28,21)(5,16,8,29,22) \\
  (6,17,9,30,23)
\end{gather*}
\subsection*{$(\alpha,\beta)=(15,2)$, lane $p=7$}
\smallskip\noindent $G_{1}$: a $C_5$-factor, \emph{base}.
\begin{gather*}
  (0,23,31,21,28)(1,10,20,15,32)(2,6,17,33,12) \\
  (3,11,8,4,25)(5,18,27,16,30)(7,19,26,9,22) \\
  (13,24,14,34,29)
\end{gather*}
\smallskip\noindent $G_{2}$: a $C_5$-factor, \emph{base}.
\begin{gather*}
  (0,17,1,20,26)(2,8,24,10,16)(3,18,9,4,21) \\
  (5,12,32,7,31)(6,28,13,14,33)(11,15,27,19,30) \\
  (22,23,29,25,34)
\end{gather*}
\smallskip\noindent $G_{3}$: a $C_5$-factor, $\rho$-invariant.
\begin{gather*}
  (0,27,8,15,30)(1,21,9,16,31)(2,22,10,17,32) \\
  (3,23,11,18,33)(4,24,12,19,34)(5,25,13,20,28) \\
  (6,26,7,14,29)
\end{gather*}
\smallskip\noindent $G_{4}$: a $C_7$-factor, $\rho$-invariant.
\begin{gather*}
  (0,1,2,3,4,5,6)(7,9,11,13,8,10,12) \\
  (14,15,16,17,18,19,20)(21,24,27,23,26,22,25) \\
  (28,31,34,30,33,29,32)
\end{gather*}
\smallskip\noindent $G_{5}$: a $C_7$-factor, $\rho$-invariant.
\begin{gather*}
  (0,2,4,6,1,3,5)(7,8,9,10,11,12,13) \\
  (14,17,20,16,19,15,18)(21,23,25,27,22,24,26) \\
  (28,29,30,31,32,33,34)
\end{gather*}
\subsection*{$(\alpha,\beta)=(13,4)$, lane $p=5$}
\smallskip\noindent $G_{1}$: a $C_5$-factor, \emph{base}.
\begin{gather*}
  (0,21,9,5,33)(1,15,16,27,26)(2,25,24,13,29) \\
  (3,18,11,10,20)(4,28,19,6,31)(7,14,23,22,17) \\
  (8,30,32,12,34)
\end{gather*}
\smallskip\noindent $G_{2}$: a $C_5$-factor, \emph{base}.
\begin{gather*}
  (0,20,29,34,26)(1,10,27,31,30)(2,18,7,25,32) \\
  (3,11,5,13,16)(4,9,24,6,22)(8,15,14,12,17) \\
  (19,23,28,21,33)
\end{gather*}
\smallskip\noindent $G_{3}$: a $C_5$-factor, $\sigma$-invariant.
\begin{gather*}
  (0,1,2,3,4)(5,10,19,32,26)(6,11,15,33,27) \\
  (7,12,16,34,28)(8,13,17,30,29)(9,14,18,31,25) \\
  (20,22,24,21,23)
\end{gather*}
\smallskip\noindent $G_{4}$: a $C_5$-factor, $\sigma$-invariant.
\begin{gather*}
  (0,2,4,1,3)(5,7,9,6,8)(10,29,21,19,31) \\
  (11,25,22,15,32)(12,26,23,16,33)(13,27,24,17,34) \\
  (14,28,20,18,30)
\end{gather*}
\smallskip\noindent $G_{5}$: a $C_5$-factor, $\sigma$-invariant.
\begin{gather*}
  (0,6,10,33,24)(1,7,11,34,20)(2,8,12,30,21) \\
  (3,9,13,31,22)(4,5,14,32,23)(15,17,19,16,18) \\
  (25,27,29,26,28)
\end{gather*}
\smallskip\noindent $G_{6}$: a $C_7$-factor, $\sigma$-invariant.
\begin{gather*}
  (0,7,26,19,22,30,11)(1,8,27,15,23,31,12) \\
  (2,9,28,16,24,32,13)(3,5,29,17,20,33,14) \\
  (4,6,25,18,21,34,10)
\end{gather*}
\smallskip\noindent $G_{7}$: a $C_7$-factor, $\sigma$-invariant.
\begin{gather*}
  (0,8,32,21,15,25,12)(1,9,33,22,16,26,13) \\
  (2,5,34,23,17,27,14)(3,6,30,24,18,28,10) \\
  (4,7,31,20,19,29,11)
\end{gather*}
\smallskip\noindent $G_{8}$: a $C_7$-factor, $\sigma$-invariant.
\begin{gather*}
  (0,9,26,11,23,33,17)(1,5,27,12,24,34,18) \\
  (2,6,28,13,20,30,19)(3,7,29,14,21,31,15) \\
  (4,8,25,10,22,32,16)
\end{gather*}
\smallskip\noindent $G_{9}$: a $C_7$-factor, $\sigma$-invariant.
\begin{gather*}
  (0,10,23,9,29,16,31)(1,11,24,5,25,17,32) \\
  (2,12,20,6,26,18,33)(3,13,21,7,27,19,34) \\
  (4,14,22,8,28,15,30)
\end{gather*}
\subsection*{$(\alpha,\beta)=(11,6)$, lane $p=5$}
\smallskip\noindent $G_{1}$: a $C_5$-factor, \emph{base}.
\begin{gather*}
  (0,1,4,16,15)(2,5,7,24,6)(3,20,18,33,23) \\
  (8,9,27,29,22)(10,11,13,34,21)(12,17,30,31,19) \\
  (14,25,26,32,28)
\end{gather*}
\smallskip\noindent $G_{2}$: a $C_5$-factor, \emph{base}.
\begin{gather*}
  (0,19,26,20,34)(1,32,2,9,33)(3,12,4,10,24) \\
  (5,11,8,18,14)(6,22,7,28,31)(13,21,23,30,25) \\
  (15,17,27,16,29)
\end{gather*}
\smallskip\noindent $G_{3}$: a $C_5$-factor, $\sigma$-invariant.
\begin{gather*}
  (0,5,16,13,28)(1,6,17,14,29)(2,7,18,10,25) \\
  (3,8,19,11,26)(4,9,15,12,27)(20,21,22,23,24) \\
  (30,32,34,31,33)
\end{gather*}
\smallskip\noindent $G_{4}$: a $C_7$-factor, $\sigma$-invariant.
\begin{gather*}
  (0,6,19,30,12,24,29)(1,7,15,31,13,20,25) \\
  (2,8,16,32,14,21,26)(3,9,17,33,10,22,27) \\
  (4,5,18,34,11,23,28)
\end{gather*}
\smallskip\noindent $G_{5}$: a $C_7$-factor, $\sigma$-invariant.
\begin{gather*}
  (0,10,8,27,24,19,33)(1,11,9,28,20,15,34) \\
  (2,12,5,29,21,16,30)(3,13,6,25,22,17,31) \\
  (4,14,7,26,23,18,32)
\end{gather*}
\smallskip\noindent $G_{6}$: a $C_7$-factor, $\sigma$-invariant.
\begin{gather*}
  (0,12,32,9,29,20,16)(1,13,33,5,25,21,17) \\
  (2,14,34,6,26,22,18)(3,10,30,7,27,23,19) \\
  (4,11,31,8,28,24,15)
\end{gather*}
\smallskip\noindent $G_{7}$: a $C_7$-factor, $\sigma$-invariant.
\begin{gather*}
  (0,18,21,32,13,8,25)(1,19,22,33,14,9,26) \\
  (2,15,23,34,10,5,27)(3,16,24,30,11,6,28) \\
  (4,17,20,31,12,7,29)
\end{gather*}
\smallskip\noindent $G_{8}$: a $C_7$-factor, $\sigma$-invariant.
\begin{gather*}
  (0,23,17,5,31,14,27)(1,24,18,6,32,10,28) \\
  (2,20,19,7,33,11,29)(3,21,15,8,34,12,25) \\
  (4,22,16,9,30,13,26)
\end{gather*}
\smallskip\noindent $G_{9}$: a $C_7$-factor, $\sigma$-invariant.
\begin{gather*}
  (0,24,14,15,6,33,26)(1,20,10,16,7,34,27) \\
  (2,21,11,17,8,30,28)(3,22,12,18,9,31,29) \\
  (4,23,13,19,5,32,25)
\end{gather*}

\section{Checking a certificate}\label{sec:check}

Fix one of the four lists. Checking it is finite and needs nothing but
the printed cycles.

\begin{enumerate}
\item Each printed object is a two-factor: its cycles are
  vertex-disjoint, they exhaust $\{0,\dots,34\}$, and all have the
  declared length, so the object has $35$ edges and the declared cycle
  type.
\item Each object marked invariant is fixed setwise by the lane map:
  applying $\rho$ (or $\sigma$) to every printed edge returns the same
  edge set.
\item Each base object $F$ has the $p$ images
  $F,\rho F,\dots,\rho^{p-1}F$ pairwise distinct, with
  $\rho^{p}F=F$ (respectively for $\sigma$); these $p$ factors have
  the same cycle type as $F$, the lane map being an automorphism of
  $K_{35}$.
\item The expansion has $17$ factors, and their $17\times35=595$ edges
  are pairwise distinct and therefore exhaust $E(K_{35})$, whose size is
  $\binom{35}{2}=595$. Equivalently, the union equals $E(K_{35})$ as a
  set.
\item Counting factors by type gives the claimed $(\alpha,\beta)$, and
  $\alpha+\beta=17$.
\end{enumerate}

Step 4 is the only laborious one. It is where the orbit bookkeeping of
Section~2 helps: the edges of an invariant factor form $35/p$ complete
edge orbits and those of a base object $35$ complete orbits, so it
suffices to name the orbits used and check that the $85$ (respectively
$119$) orbits are each named once. A machine check of all five steps for
all four cells, in exact integer arithmetic, accompanies this note
(\texttt{verify.py}, $66$ checks, all passing); it reads the cycle lists
exactly as printed above. It checks the five steps only: the hypothesis
of Corollary~\ref{cor:complete} is quoted to it as a list of fourteen
values of $\beta$, so the program tests the arithmetic of assembling the
spectrum and not that hypothesis.

\section{Scope}

Only the instance $(m,n)=(5,7)$, $v=mn=35$, is treated; the other cells
of the exception clause of \cite{BDT17} are untouched. The cell $(9,8)$
is not ours: it is due to \cite{WLC}, and it is not printed here. We
make no claim about uniqueness, enumeration, classification, minimality
or nonexistence.

We make no priority claim either. What the four objects establish is
that the four pairs are realizable, which is a statement about the
pairs and not about the literature: it is not asserted here that their
realizability was previously unrecorded, and no literature check
settling that is reported. Falling outside the positive range of the two
theorems cited above is a fact about those two theorems, not evidence
that no other source realizes these cells.

\begin{thebibliography}{9}

\bibitem{BDT17}
A.~C.~Burgess, P.~Danziger and T.~Traetta,
\emph{On the Hamilton--Waterloo problem with odd orders},
J. Combin. Des. \textbf{25} (2017), 258--287.
\href{https://doi.org/10.1002/jcd.21552}{doi:10.1002/jcd.21552};
\href{https://arxiv.org/abs/1510.07079}{arXiv:1510.07079}

\bibitem{WLC}
L.~Wang, S.~Lu and H.~Cao,
\emph{Further results on almost resolvable cycle systems and the
Hamilton--Waterloo problem},
J. Combin. Des. \textbf{26} (2018), 27--47.
\href{https://doi.org/10.1002/jcd.21571}{doi:10.1002/jcd.21571};
\href{https://arxiv.org/abs/1605.00818}{arXiv:1605.00818}

\end{thebibliography}

\end{document}
